\documentclass[10pt,a4paper]{amsart}
\usepackage[margin=19mm]{geometry}
\usepackage{amsmath,amssymb,mathtools}
\usepackage[scr=rsfs,cal=euler]{mathalfa}
\usepackage{xfrac}
\usepackage[unicode,hidelinks,bookmarks=false]{hyperref}
\newcommand{\ie}{i.e.\ }
\newcommand{\cf}{cf.\ }
\newcommand{\resp}{respectively\ }
\newcommand{\Subsection}{\subsection}
\providecommand{\nobreakdash}{\nobreak\hbox{-}}
\setlength{\emergencystretch}{2em}
\multlinegap0pt
\usepackage{bm}
\usepackage{bbm}
\usepackage{dsfont}
\usepackage{xspace}
\usepackage{cancel}
\usepackage{mathtools}
\usepackage{amsthm}
\makeatletter
\@ifundefined{definition}{\newtheorem{definition}{Definition}}{}
\@ifundefined{lemma}{\newtheorem{lemma}[definition]{Lemma}}{}
\makeatother
\title[Higher-order spectral perturbation expansions II: Kernel mat]{Higher-order spectral perturbation expansions II: Kernel matrices and manifold learning}
\author{Bernhard Stankewitz, Martin Wahl}
\date{Source: arXiv:2606.16373v1 (CC BY 4.0)}
\begin{document}
\maketitle

\noindent\emph{Source and licence.} Higher-order spectral perturbation expansions II: Kernel matrices and manifold learning. Version arXiv:2606.16373v1 (CC BY 4.0), distributed under the Creative Commons Attribution 4.0 International licence, as recorded in the arXiv licence metadata for this exact version and shown on the abstract page https://arxiv.org/abs/2606.16373v1. Full rights notice: licenses/arxiv-2606.16373-CC-BY-4.0.txt.\par\smallskip

\noindent\textit{Editorial scope.} Counted unit: the Lemma whose source label is \texttt{lem\_TaylorRemainder} in the article (source0.tex lines 1462--1492, source label \texttt{lem\_TaylorRemainder}), delivered together with the article's own complete original proof of it (source0.tex lines 1494--1764, one \texttt{proof} environment of 271 lines). No mathematics is written, repaired or re-derived: statement and proof are the article's own text line for line. Required context retained uncounted: lem\_EigenvalueSeparation (lines 1362--1384). Every definition, display and label of those lines is retained unchanged. Omitted from this package and not redistributed: 1--1361, 1385--1461, 1493--1493, 1765--2906. Nothing inside a retained excerpt is omitted or altered: every retained body is delivered exactly as the article's lines are written, so no omission receipt of any kind accompanies this package. Citations: the retained excerpts carry no citation command at all, so no bibliography entry of the article is transcribed and no reference is redistributed. Reference adaptation: none was needed and none was made; every reference inside the delivered unit resolves inside it, so no unresolved reference or citation is produced. Printed slips kept literal: no printed word, symbol, hypothesis or number of the counted statement or of its proof was altered. Layout and class: the article's own class is \texttt{article} with packages including fullpage, babel, inputenc, biblatex, authblk, amssymb, mathtools, amsthm, amsmath, upgreek, mathrsfs, graphicx, float, enumitem; the wrapper is the lane's pinned \texttt{amsart} editorial wrapper at 10pt on A4, which restates the theorem-like environments the retained text needs and adds the article's own macro definitions that the retained text needs (none), with extra wrapper packages (bm, bbm, dsfont, xspace, cancel, mathtools). The delivered numbering is the unit's own. Assets: no retained span contains a figure environment, a graphics-inclusion command, a PDF-page inclusion, a table or a TikZ picture; no image file is copied into this package, the package declares no asset dependency and no image file is redistributed with it. Licence: version arXiv:2606.16373v1 of this article is distributed under the Creative Commons Attribution 4.0 International licence, as recorded in the arXiv licence metadata for this exact version and shown on the abstract page https://arxiv.org/abs/2606.16373v1. A full rights notice is delivered at licenses/arxiv-2606.16373-CC-BY-4.0.txt. Characters: every retained excerpt is pure ASCII, so the delivered engine, XeLaTeX with its default Latin Modern text font, reports no missing character.
\begin{lemma}[Eigenvalue separation]%{{{
  \label{lem_EigenvalueSeparation}
  If \( \delta_{r} < 1 \), then we have 
  \begin{align}
            | \widehat{ \lambda }_{j} - \mu_{r} | 
    & \le
            \delta_{r} g_{r} 
            \qquad \qquad \qquad
            \qquad \text{for all } j \in I_{r},
    \\
            \widehat{ \lambda }_{k} - \mu_{r + 1} 
    & \le 
            \delta_{r} ( \mu_{r} - \mu_{r + 1} ) 
            \quad \qquad \ \ \text{for all } k \in I_{r + 1},
    \notag
    \\
            \widehat{ \lambda }_{k} - \mu_{r - 1} 
    & \ge
            - \delta_{r} ( \mu_{r - 1} - \mu_{r} ) 
            \qquad \ \ \; \text{for all } k \in I_{r - 1}.
    \notag
  \end{align}
\end{lemma}%}}}

\begin{lemma}[Taylor remainder]%{{{
  \label{lem_TaylorRemainder}
  Suppose that \( \delta_{r} < 1 \).
  Then, for any integer \( p \ge 1 \), 
  \begin{align}
    \label{eq_A_TaylorRemainder}
        \widehat{P}_{r} - \sum_{n = 0}^{p - 1} P_{r}^{ (n) }
    & = 
        (-1)^{p - 1} 
        \sum_{k_{1}, \dots k_{p}, \ge 0} 
        R_{r}^{ ( k_{1} ) } E \dots E R_{r}^{ ( k_{p} ) } 
        E \widehat{R}_{r}^{ ( p - k_{1} - \dots - k_{p} ) }
  \end{align}
  with 
  \( 
    \widehat{R}_{r}^{ (k) } = - \sum_{ j \in I_{r} }  
                                ( \widehat{ \lambda }_{j} - \mu_{r} )^{ - k }
                                \widehat{u}_{j} \otimes \widehat{u}_{j} 
  \),
  for \( k \le 0 \) and \( \widehat{R}_{r}^{ (k) } = \widehat{R}_{r} ^{k} \),
  for  \( k > 0 \), where 
  \begin{align}
    \widehat{R}_{r}: = \sum_{ s \ne r } \sum_{ j \in I_{s} } 
                       ( \widehat{ \lambda }_{j} - \mu_{r}  )^{-1} 
                       \widehat{u}_{j} \otimes \widehat{u}_{j}.
  \end{align}
  is the empirical reduced resolvent.
  It is well defined, since \( \widehat{ \lambda }_{j} - \mu_{r} \ne 0 \) for 
  \( j \in I_{s}, s \ne r \)  due to Lemma 
  \normalfont \ref{lem_EigenvalueSeparation}.
\end{lemma}%}}}

\begin{proof} %{{{
  We prove the statement in Equation \eqref{eq_A_TaylorRemainder} via induction
  over \( p \in \mathbb{N} \).
  For \( p = 1 \), we write \( \widehat{P}_{r} - P_r \) as
  \begin{align}
        \widehat{P}_{r} - P_{r} 
    & = 
        ( I - P_{r} ) \widehat{P}_{r} - P_{r} ( I - \widehat{P}_{r} ).
  \end{align}
  Using \begin{align}
  \label{eq_ProjectureToEigenvalueEquality}
  ( u_{j} \otimes u_{j} ) E ( \widehat{u}_{k} \otimes \widehat{u}_{k} ) 
  & = 
  ( \widehat{ \lambda }_{k} - \lambda_{j} ) 
  ( u_{j} \otimes u_{j} ) ( \widehat{u}_{k} \otimes \widehat{u}_{k} ) 
\end{align}
  and Lemma \ref{lem_EigenvalueSeparation}, the second term on the right-hand is given by
  \begin{align}
      P_{r} ( I - \widehat{P}_{r} ) 
    = 
      \sum_{ j \in I_{r} } 
      \sum_{ s \ne r } \sum_{ k \in I_{s} }
      ( \widehat{ \lambda }_{k} - \mu_{r} )^{-1}
      ( u_{j} \otimes u_{j} ) E ( \widehat{u}_{k} \otimes \widehat{u}_{k} ) 
    = 
      P_{r} E \widehat{R}_{r}.
  \end{align}
  Further, with analogous reasoning, the first term on the right-hand side is 
  given by
  \begin{align}
      ( I - P_{r} ) \widehat{P}_{r} 
    =
      \sum_{ k \in I_{r} } 
      \sum_{ s \ne r } ( \widehat{ \lambda }_{k} - \mu_{s} )^{-1} P_{s}
      E ( \widehat{u}_{k} \otimes \widehat{u}_{k} ) 
    = 
      - \sum_{ k \in I_{r} } \sum_{ l = 1 }^{ \infty } 
      ( \widehat{ \lambda }_{k} - \mu_{r} )^{ l - 1 } 
      R_{r}^{l} E ( \widehat{u}_{k} \otimes \widehat{u}_{k} ). 
    \notag
  \end{align}
  For the last equality, we have used a geometric series expansion of
  \( 
      ( \widehat{ \lambda }_{k} - \mu_{s} )^{-1} 
    = 
      - ( \mu_{s} - \mu_{r} )^{-1} / 
        ( 1 - ( \widehat{ \lambda }_{k} - \mu_{r} ) / ( \mu_{s} - \mu_{r} ) )
  \),
  which is well defined due to Lemma \ref{lem_EigenvalueSeparation}.
  It follows that 
  \begin{align}
            \widehat{P}_{r} - P_{r} 
    & = 
            \sum_{ l = 1 }^{ \infty } 
            R_{r}^{l} E \widehat{R}_{r}^{ ( 1 - l ) }
          + 
            R_{r}^{ (0) } E \widehat{R}_{r}^{ (1) }
    = 
            \sum_{ l = 0 }^{ \infty } 
            R_{r}^{l} E \widehat{R}_{r}^{ ( 1 - l ) },
    \notag
  \end{align}
  which establishes the induction hypothesis.

  For the induction step, we assume that the statement is true for a fixed 
  \( p \) and consider the fact that the induction hypothesis can be written as
  \begin{align}
          \widehat{R}_{r}^{ (0) } 
    & =
          R_{r}^{ (0) } 
        - 
          \sum_{l = 0}^{ \infty } 
          R_{r}^{ (l) } E \widehat{R}_{r}^{ (1 - l) }.
  \end{align}
  We will show that analogously
  \begin{align}
    \label{eq_RKSmallerZero}
          \widehat{R}_{r}^{ (k) } 
    & = 
        \qquad \
        - \sum_{l = 0}^{ \infty } 
          R_{r}^{ (l) } E \widehat{R}_{r}^{ (k + 1 - l) },
        \qquad k < 0;
    \\
    \label{eq_RKLargerZero}
          \widehat{R}_{r}^{ (k) } 
    & = 
          R_{r}^{ (k) } 
        - 
          \sum_{l = 0}^{ \infty } 
          R_{r}^{ (l) } E \widehat{R}_{r}^{ (k + 1 - l) }, 
        \qquad k \ge 0.
  \end{align}
  Setting \( k: = p - k_{1} - \dots - k_{p} \), it then follows from Equations
  \eqref{eq_RKSmallerZero}, \eqref{eq_RKLargerZero} and the induction hypothesis
  that 
  \begin{align}
         \widehat{P}_{r} - \sum_{n = 0}^{p - 1} P_{r}^{ (n) } 
    & =  
         (-1)^{p - 1} \sum_{k < 0} 
         \sum_{ \substack{ k_{1}, \dots, k_{p} \ge 0 \\
                           k_{1} + \dots + k_{p} = p - k 
                         }
              }
         R_{r}^{ ( k_{1} ) } E \dots E R_{r}^{ ( k_{p} ) } E
         \Big( - \sum_{l = 0}^{ \infty } 
                 R_{r}^{ (l) } E \widehat{R}_{r}^{ (k + 1 - l) } 
         \Big)
    \notag
    \\
    & + 
            (-1)^{p - 1} \sum_{k \ge 0} 
            \sum_{ \substack{ k_{1}, \dots, k_{p} \ge 0 \\
                              k_{1} + \dots + k_{p} = p - k 
                            }
              }
            R_{r}^{ ( k_{1} ) } E \dots E R_{r}^{ ( k_{p} ) } E
            \Big( 
              R_{r}^{ (k) } 
            - 
              \sum_{l = 0}^{ \infty } 
              R_{r}^{ (l) } E \widehat{R}_{r}^{ ( k + 1 - l ) } 
            \Big)
    \notag
    \\
    & = 
            (-1)^{p} 
            \sum_{k_{1}, \dots, k_{p + 1} \ge 0} 
            R_{r}^{ ( k_{1} ) } E \dots E R_{r}^{ ( k_{p + 1} ) } E 
            \widehat{R}_{r}^{ (p + 1 - k_{1} - \dots - k_{p + 1}) } 
            +
            P_{r}^{ (p) },
    \notag
  \end{align}
  which yields the statement in Lemma \ref{lem_TaylorRemainder} by rearranging
  the terms.

  It remains to be shown that Equation \eqref{eq_RKSmallerZero} and
  \eqref{eq_RKLargerZero} are true for \( k < 0 \) and \( k > 0 \) respectively. 
  For \( k < 0 \), we have that
  \begin{align}
          \widehat{R}_{r}^{ (k) } 
    & = 
        - \sum_{ j \in I_{r} } 
          ( \widehat{ \lambda }_{j} - \mu_{r} )^{-k} 
          P_{r} ( \widehat{u}_{j} \otimes \widehat{u}_{j} )
        - \sum_{ j \in I_{r} } 
          ( \widehat{ \lambda }_{j} - \mu_{r} )^{-k} 
          ( I - P_{r} ) ( \widehat{u}_{j} \otimes \widehat{u}_{j} ),
    \notag 
  \end{align}
  where, using Equation \eqref{eq_ProjectureToEigenvalueEquality}, the first
  term on the right-hand side is equal to
  \begin{align}
      - \sum_{ j \in I_{r} } 
        ( \widehat{ \lambda }_{j} - \mu_{r} )^{ - k - 1 } 
        P_{r} E ( \widehat{u}_{j} \otimes \widehat{u}_{j} ) 
    = 
      - R_{r}^{ (0) } E \widehat{R}_{r}^{ ( k + 1 ) }.
  \end{align}
  Reusing the geometric series expansion from before, the second term on the
  right-hand side can be written as
  \begin{align}
    & \ \ \ \
        - \sum_{s \ne r} \sum_{ j \in I_{r} } \sum_{ j' \in I_{s} } 
          ( \widehat{ \lambda }_{j} - \mu_{r} )^{-k} 
          ( \widehat{ \lambda }_{j} - \mu_{s} )^{-1}
          ( u_{j'} \otimes u_{j'} ) E ( \widehat{u}_{j} \otimes \widehat{u}_{j} ) 
    \\
    & = 
        - \sum_{ j \in I_{r} } ( \widehat{ \lambda }_{j} - \mu_{r} )^{-k} 
          \sum_{ s \ne r } 
          ( \widehat{ \lambda }_{j} - \mu_{s} )^{-1}
          P_{s} E ( \widehat{u}_{j} \otimes \widehat{u}_{j} )
    \notag
    \\
    & =
          \sum_{ j \in I_{r} } \sum_{l = 1}^{ \infty } 
          ( \widehat{ \lambda }_{j} - \mu_{r} )^{-k - 1 + l} 
          R_{r}^{ (l) } E ( \widehat{u}_{j} \otimes \widehat{u}_{j} ) 
    = 
        - \sum_{l = 1}^{ \infty } 
          R_{r}^{ (l) } E \widehat{R}_{r}^{ (k + 1 - l) }.
    \notag
  \end{align}
  Together, this yields Equation \eqref{eq_RKSmallerZero}.

  For \( k > 0 \), we write
  \begin{align}
    \label{eq_RKLargerZeroExpanded_1}
        \widehat{R}_{r}^{ (k) } 
    & = 
        ( I - P_{r} ) \widehat{R}_{r}^{k} + P_{r} \widehat{R}_{r}^{k}.
  \end{align}
  Focusing on the first term on the right-hand side, we note that inserting \eqref{eq_ProjectureToEigenvalueEquality} into \( R_{r} E \widehat{R}_{r} \) yields \( ( I - P_{r} ) \widehat{R}_{r} = R_{r} ( I - \widehat{P}_{r} ) - R_{r} E \widehat{R}_{r} \)
  and inductively
  \begin{align}
    \label{eq_RKLargerZeroExpanded_2}
            ( I - P_{r} ) \widehat{R}_{r}^{k} 
    & = 
            R_{r}^{k} ( I - \widehat{P}_{r} ) 
          - 
            \sum_{ l = 1 }^{k} 
            R_{r}^{l} E \widehat{R}_{r}^{k + 1 - l}.
  \end{align}
  Together, Equations \eqref{eq_RKLargerZeroExpanded_1} and
                      \eqref{eq_RKLargerZeroExpanded_2} then amount to
  \begin{align}
            \widehat{R}_{r}^{k} 
    & = 
            R_{r}^{k} ( I - \widehat{P}_{r} ) 
          - 
            \sum_{ l = 1 }^{k} 
            R_{r}^{l} E \widehat{R}_{r}^{k + 1 - l}
          + 
            P_{r} \widehat{R}_{r}^{k}
    \\
    & = 
            R_{r}^{k} 
          - 
            \sum_{ l = k + 1 }^{ \infty } 
            R_{r}^{l} E \widehat{R}_{r}^{(l - 1 - k)} 
          - 
            \sum_{ l = 1 }^{k} 
            R_{r}^{l} E \widehat{R}_{r}^{k + 1 - l}
          + 
            P_{r} \widehat{R}_{r}^{k}
    \notag
    \\
    & = 
    \notag
        R_{r}^{k} 
      - 
        \sum_{ l = 1 }^{ \infty } 
        R_{r}^{l} E \widehat{R}_{r}^{ (l - 1 - k) } 
      + 
        P_{r} E \widehat{R}_{r}^{k + 1}
    = 
        R_{r}^{k} 
      - 
        \sum_{l = 0}^{ \infty } 
        R_{r}^{l} E \widehat{R}_{r}^{ (l - 1 - k) },
  \end{align}
  where again, we have used Equation \eqref{eq_ProjectureToEigenvalueEquality}
  and the geometric series expansion to obtain
  \begin{align}
        - R_{r}^{k} \widehat{P}_{r} 
    & = 
        - R_{r}^{k} ( I - P_{r} ) \widehat{P}_{r} 
    = 
        - R_{r}^{k} 
          \sum_{ s \ne r } 
          \sum_{ j  \in I_{s} }
          \sum_{ j' \in I_{r} } 
          ( \widehat{ \lambda }_{j'} - \lambda_{j} )^{-1} 
          ( u_{j} \otimes u_{j} ) E 
          ( \widehat{u}_{j'} \otimes \widehat{u}_{j'}) 
    \\
    & = 
          R_{r}^{k}
          \sum_{ j' \in I_{r} } 
          \sum_{ l = 1 }^{ \infty } 
          ( \widehat{ \lambda }_{ j' } - \mu_{r} )^{ l - 1 } R_{r}^{l}
          E ( \widehat{u}_{ j' } \otimes \widehat{u}_{ j' }) 
    = 
        - \sum_{ l = k + 1 }^{ \infty } 
          R_{r}^{l} E \widehat{R}_{r}^{ ( l - 1 - k ) }.
    \notag
  \end{align}
  This finishes the proof.
\end{proof}%}}}

\end{document}
